\documentclass[10pt,a4paper]{amsart}
\usepackage[margin=19mm]{geometry}
\usepackage{amsmath,amssymb,mathtools}
\usepackage[scr=rsfs,cal=euler]{mathalfa}
\usepackage{xfrac}
\usepackage[unicode,hidelinks,bookmarks=false]{hyperref}
\newcommand{\ie}{i.e.\ }
\newcommand{\cf}{cf.\ }
\newcommand{\resp}{respectively\ }
\newcommand{\Subsection}{\subsection}
\providecommand{\nobreakdash}{\nobreak\hbox{-}}
\setlength{\emergencystretch}{2em}
\multlinegap0pt
\usepackage{amssymb}
\usepackage{dsfont}
\usepackage[shortlabels]{enumitem}
\usepackage{tikz}
\usepackage{tcolorbox}
\newtheorem{theorem}{Theorem}
\DeclareMathOperator{\R}{\mathbb{R}}
\renewcommand{\d}{\mathrm{d}}
\DeclareMathOperator{\pv}{\operatorname{p.\!v.}}
\title{Optimal stability of complement value problems for $p$-L{\'e}vy operators.}
\author{Guy Foghem}
\date{Source: arXiv:2605.13389v1 (CC BY 4.0)}
\begin{document}
\maketitle

\noindent\emph{Source and licence.} Optimal stability of complement value problems for $p$-L{\'e}vy operators.. Version arXiv:2605.13389v1 (CC BY 4.0), distributed by arXiv under the Creative Commons Attribution 4.0 International licence, http://creativecommons.org/licenses/by/4.0/, as recorded in the arXiv licence metadata for this exact version and shown on the abstract page https://arxiv.org/abs/2605.13389v1. Full licence notice: \texttt{licenses/arxiv-2605.13389-CC-BY-4.0.txt}.\par\smallskip

\noindent\textit{Editorial scope.} Counted unit: the article's theorem thm:point-evalua-pLevy-rad (source0.tex lines 1547--1575). The article's complete original proof of it is delivered with it (source0.tex lines 1577--1636, one \texttt{proof} environment of 60 lines). No mathematics is written, repaired or re-derived: the statement and the proof are the article's own lines, in the article's own notation, and no retained line is substituted or re-flowed.  No further source-local context is needed: every cross-reference of every retained line resolves inside the two delivered spans.  Nothing was omitted from any retained span, so the package carries no omission record.  No retained line contains a citation, so the wrapper carries no bibliography and no bibliography entry.  No retained line contains a figure environment, an image-inclusion command or drawing code, so no image is delivered and the package declares no asset dependency.  Editorial formatting only, in addition: an AMS \texttt{amsart} preamble that restates the article's own declarations and macros used by the retained lines, namely the environments \texttt{theorem} on separate counters, together with the standard packages those lines need.  The retained environments print in this document as Theorem 1 in order of appearance, whereas the article prints the same items with the numbering of its own full document.  The complete native source of the article as distributed in the arXiv e-print for this exact version is bundled unchanged as \texttt{sources/200498/source0.tex}.
\begin{theorem}\label{thm:point-evalua-pLevy-rad}
Let $\nu \in L^1(\R^d,\, 1 \wedge |h|^p)$ be a non-negative $p$-L\'evy radial kernel, and whose corresponding $p$-L\'evy operator is given by
\begin{align*}
Lu(x)= \pv 2\int_{\R^d}|u(x)-u(y)|^{p-2}(u(x)-u(y))\, \nu(x-y)\d y.
\end{align*}
Let $u \in C^{2}(B_{2}(x)) \cap L^{\infty}(\R^d)$ with $x \in \R^d$. Assume that $|\nabla u(x)| \neq 0$ only when $1 < p < 2$.
Then $Lu(x)$ is well defined, and for  $0<\delta<1$ small
\begin{align*}
|Lu(x)|&\leq 2^p\|u\|^{p-1}_{L^\infty(\R^d)}\int_{|h|\geq \delta}\nu(h)\d h\\
+&2^{p+1}p\Big(\|u\|_{C^2(B_\delta(x))} \frac{K_{d,p-2}}{|\nabla u(x)|^{2-p}}+
\|u\|^{\max(1, p-1)}_{C^2(B_\delta(x))}\Big)\int_{|h|<\delta}|h|^p\, \nu(h)\d h
\end{align*}
with  $K_{d,p-2}= \tfrac{d+p-2}{p-1} K_{d,p}. $ In particular, for some constant $c_{d,p}>0$ we have
\begin{align*}
\begin{split}
|Lu(x)|\leq c_{d,p}\Big(\|u\|_{C^2(B_1(x))} |\nabla u(x)|^{p-2}+
\|u\|^{\max(1, p-1)}_{C^2(B_1(x))}+\|u\|^{p-1}_{L^\infty(\R^d)}\Big)
\|\nu\|_{L^1(\R^d,1\land|h|^p)}.
\end{split}
\end{align*}
Furthermore,  if $u \in C^{2}_b(U) \cap L^{\infty}(\mathbb{R}^d)$ where $U\subset  \R^d$ is open, then for $p\geq 2$ we have
\begin{align*}
\|Lu\|_{L^\infty(U)}\leq c_{d,p}\Big(\|u\|^{p-1}_{C^2_b(U)} + \|u\|^{p-1}_{L^\infty(\R^d)}\Big)\|\nu\|_{L^1(\R^d,1\land|h|^p)}.
\end{align*}
 Meanwhile, for the case $1<p<2$,  if in addition $\frac{1}{|\nabla u|}\in L^\infty(U)$  then  we have
\begin{align*}
\|Lu\|_{L^\infty(U)}\leq c_{d,p} \Big( \|\, |\nabla u|^{-1}\|^{2-p}_{L^\infty(U)}\|u\|_{C^2_b(U)}  +\|u\|_{C^2_b(U)}  + \|u\|^{p-1}_{L^\infty(\R^d)}\Big)\|\nu\|_{L^1(\R^d,1\land|h|^p)}.
\end{align*}
\end{theorem}

\begin{proof}
We prove only the first estimate, as the remaining ones follow directly from it.   Let us put $a=\nabla u(x)\cdot h$ and $b=\int_0^1\nabla u(x+th)\cdot h\, \d t$ so that $u(x+h)-u(x)= \int_0^1\nabla u(x+th)\cdot h\d t$. Since  $\pv \int_{B_\delta(0)} \psi(\nabla u(x)\cdot h)\nu(h)\d h=0$ and $\psi(t)= |t|^{p-2}t$, by the fundamental theorem of calculus we obtain
\begin{align}\label{eq:xxpsi-taylor-nearzero}
\begin{split}
2\int_{|h|<\delta} &\psi(u(x+h)-u(x))\,\nu(h)\,\d h
\\&= 2\int_{|h|<\delta } \psi'(a)(b-a)+ [\psi(b)-\psi(a)-\psi'(a)(b-a)]\,\nu(h)\,\d h
\end{split}
\end{align}
where we emphasize that $\psi'(t)= (p-1)|t|^{p-2}$. It is worth noticing that
\begin{align*}
2|b-a| &=2\Big| \int_0^1 t\int_0^1 D^2u(x+sth)\cdot h \cdot h\, \d s\,  \d t\Big|\leq \|u\|_{C^2(B_1(x))} |h|^2.
\end{align*}
While, since   $\nabla u(x)\neq 0$ for $1<p<2$, using polar coordinates we find that
\begin{align*}
 2\Big|\int_{|h|<\delta } \psi'(a)&(b-a)\,\nu(h)\,\d h\Big|
\leq (p-1)\|u\|_{C^2(B_1(x))}\int_{|h|<\delta }  |\nabla u(x)\cdot h|^{p-2}  |h|^2\nu(h)\,\d h\\
&= (p-1)\|u\|_{C^2(B_1(x))}\int_{\mathbb{S}^{d-1}}  \hspace*{-2ex}|\nabla u(x)\cdot \xi|^{p-2}  \d\sigma_{d-1}(\xi)\, \int_0^\delta r^{p+d-1}\nu(r)\,\d r\\
&= (p-1)\|u\|_{C^2(B_1(x))} |\nabla u(x)|^{p-2}K_{d,p-2} \int_{|h|<\delta } |h|^p\nu(h)\,\d h,
\end{align*}
where we recall that
\begin{align*}
K_{d,p-2}= \frac{1}{|\mathbb{S}^{d-1}|} \int_{\mathbb{S}^{d-1}} |w_d|^{p-2}\d \sigma_{d-1}(w)= \frac{ \Gamma\big(\frac{d}{2}\big)  \Gamma\big(\frac{p-1}{2}\big)}{\Gamma\big(\frac{d+p-2}{2}\big) \Gamma\big(\frac{1}{2}\big)}= \frac{d+p-2}{p-1} K_{d,p}.
\end{align*}
Next, for  $1<p<2$ with $\nabla u(x)\neq 0$, note that  $\psi(t)= |t|^{p-2}t $ is $C^1$ on $\R\setminus\{0\}$  so that
\begin{align*}
\lim_{b\to a}\frac{\psi(b)-\psi(a) -\psi'(a)(b-a)}{b-a}=0.
\end{align*}
Since $b- a\to0$ as $|h|\to 0$, for $0<\delta<1$ small enough we have
\begin{align*}
|\psi(b)-\psi(a) -\psi'(a)(b-a)| \leq |b-a|\leq \|u\|_{C^2(B_1(x))}|h|^2\leq \|u\|_{C^2(B_1(x))}|h|^p.
\end{align*}
For the case $p\geq 2$ we have
\begin{align*}
|\psi(b)-\psi(a)-\psi'(a)(b-a)|
&\leq |b-a| \int_0^1  |\psi'(a+\tau (b-a))- \psi'(a)|\\
&\leq 2^{p+1}(p-1)\|u\|_{C^2(B_1(x))}
\|\nabla u\|^{p-2}_{L^\infty(B_1(x))}
|h|^p\\
&\leq 2^{p+1}(p-1)\|u\|^{p-1}_{C^2(B_1(x))} |h|^p.
\end{align*}
Thus in both cases $1<p<2$  and $p\geq 2$ we have
\begin{align*}
\footnotesize
2\Big|\int_{|h|<\delta }\hspace*{-2ex}|\psi(b)-\psi(a) -\psi'(a)(b-a)| \nu(h)
\d h\Big| \leq 2^{p+1}(p-1) \|u\|^{\max(1, p-1)}_{C^2(B_1(x))}\int_{|h|<\delta } \hspace*{-2ex}|h|^p\nu(h)\d h.
\end{align*}
Inserting altogether in \eqref{eq:xxpsi-taylor-nearzero} for $0<\delta<1$ sufficiently small we obtain
\begin{align*}
&2\Big|\int_{|h|<\delta } \psi(u(x+h)-u(x))\,\nu(h)\,\d h\Big|
\\
&\leq 2^{p+1}(p-1)
\Big( \|u\|_{C^2(B_1(x))} \frac{K_{d,p-2}}{|\nabla u(x)|^{2-p}} +
\|u\|^{\max(1, p-1)}_{C^2(B_1(x))}\Big)\int_{|h|<\delta }\hspace*{-2ex} |h|^p\nu(h)\d h.
\end{align*}
The desired estimate follows since for  $0<\delta<1$ we have
\begin{align*}
\int_{B^c_\delta(x)} |u(x)- u(y)|^{p-1}\nu(x-y)\d y
&= 2^{p}\|u\|^{p-1}_{L^\infty(\R^d)}\int_{|h|\geq \delta}  \nu(h) \d h.
\end{align*}
\end{proof}

\end{document}
